\margin{$\unsure{6|?}$}
2$^\circ$) $(X,\Sigma,S)$, avec
\[
\left|
\begin{aligned}
X&\ \text{compacte à bord orienté}\\
\Sigma&\ \text{sous-variété }\deleted{\text{\unsure{fermée}}}
 \inserted{\text{\unsure{pas compacte}}}\ \text{de dim }1\\
&\ \text{de }X\setminus\partial X\\
S\subset\Sigma\cup\partial X&\ \text{partie finie}
\end{aligned}
\right.
\]
morphismes : classes d'isotopie d'homéom. \unsure{respectant}\\
l'orientation. Notation $\mathcal{T}\unsure{\deleted{S}D|D}_{\mathrm{top}}$ (D initiale de\\
\inserted{divisé}\\
\deleted{``\unsure{sauvage}''}) : ``\uline{groupoïde de Teichmüller orienté}\\
\uline{divisé topologique}''.

\margin{\inserted{imm. ouverte}\\
(\deleted{\unsure{inclusion pleine}})}

\uline{NB}\quad
\diagram{\begin{tikzcd}[column sep=small,row sep=small]
S\mathcal{T}_{\mathrm{top}}\arrow[d,shorten <=3pt,shorten >=3pt]
  & \subset & S\mathcal{T}D_{\mathrm{top}}\\
\makebox[0pt][l]{\(\text{sous-groupoïde des objets tels que }\Sigma=\varnothing\)}
  & &
\end{tikzcd}}

3$^\circ$) On définit de $=$ $\mathcal{T}S^{\pm}_{\mathrm{top}}$
(\uline{groupoïde de}\\
\uline{Teichmüller biorienté divisé topologique})
